\textbf{Problem 2 answer} \GradescopeComment{This should be on page 5 of your PDF.}

\begin{enumerate}

\item[(a)]   Complete the following table to show,  for each $d$,
  the minimum input $N$ for which greedy with denominations $\{1,5,d\}$
  uses too many coins.  For the $d$ for which greedy works, put $\infty$.

  \BLUE{
    \[
      \begin{array}{c|c}
        d & N \\ \hline
        6 & \REPLACEME \\
        7 & \REPLACEME \\
        8 & \REPLACEME \\
        9 & \REPLACEME \\
        11 & \REPLACEME
      \end{array}
    \]
  }

\end{enumerate}

\continued

\paragraph{Problem 2(b) answer} \GradescopeComment{This should be on page 6 (top) of your PDF.}

\medskip

\noindent (b) The $d$ for which the greedy algorithm works is \BLUE{
  \(
  d=\REPLACEME
  \)
}.

\medskip

Here is the requested lemma and proof (for that $d$):

% COMPLETE THE LEMMA AND PROOF BELOW.  

\begin{lemma}
  For any \(N\ge
  \BLUE{
    \REPLACEME              % replace the \REPLACEME with the d you chose 
  }
  \),
  some optimal way of making change for $N$
  (using coin values in \(\{1, 5,
  \BLUE{
    \REPLACEME              % replace the \REPLACEME with the d you chose 
  }
  \}\))
  contains at least one \BLUE{
    \(
    \REPLACEME              % replace the \REPLACEME with the d you chose 
    \)
  }
\end{lemma}

\begin{longFormProof}

  \step Consider any such $N$.

  \step Let $X^*$ be an optimal way of making change for $N$ using coins of those values.

  \smallskip 

  \begin{block}
    \step \underline{Case 1.}  First consider the case that $X^*$ contains a \BLUE{
      \REPLACEME              % replace the \REPLACEME with the d you chose 
    }.
    \step Then some optimal way of making change for $N$ (namely $X^*$) contains a \BLUE{
      \REPLACEME              % replace the \REPLACEME with the d you chose 
    }.
  \end{block}

  \smallskip 

  \begin{block}

    \step \underline{Case 2.}  Next consider the case that \BLUE{
      \REPLACEME
    }

    \begin{blue}

      \step \REPLACEME          % your reasoning goes here

    \end{blue}

    \step So some optimal way of making change for $N$ contains a \BLUE{
      \REPLACEME              % replace the \REPLACEME with the d you chose 
    }.
    
  \end{block}

  \smallskip 

  \begin{block}

    \step \underline{Case 3.}  Otherwise \BLUE{
      \REPLACEME.
    }

    \begin{blue}
      \step \REPLACEME

      \step \REPLACEME

      \step $~~\vdots$              % your proof may need more than three steps and/or more than 3 cases

    \end{blue}

  \end{block}

  \smallskip 

  \step By Cases 1--3, some optimal way of making change for $N$ contains a
  \BLUE{
    \REPLACEME              % replace the \REPLACEME with the d you chose
  }.

\end{longFormProof}


\newpage


\paragraph{Problem \arabic{problem} collaboration and citations}~\GradescopeComment{This should be on page 7 (top) of your PDF.}

I collaborated on this problem with the following people:
\begin{blue}
  \begin{itemize}
  \item \REPLACEME                    % write "none" if none
  \end{itemize}
\end{blue}

My answers use ideas from the following sources (other than the lecture notes and textbooks): 
\begin{blue}
  \begin{itemize}
  \item \REPLACEME                    % write "none" if none
  \end{itemize}
\end{blue}



%%% Local Variables:
%%% mode: latex
%%% TeX-master: "homework_4"
%%% End:
